\section{Introduction}

Wage rigidity is a central feature of labor markets because it governs how shocks translate into wages, employment, and workers’ income risk. A large literature studies its origins, emphasizing either contractual frictions that lead firms to smooth wages over time or institutional constraints that limit wage adjustment. In standard frameworks, these forces imply a substitution between wages and employment: when wages cannot adjust, firms respond by adjusting employment instead. This paper proposes a different perspective. Wage rigidity can arise endogenously from firms’ selection of which workers to retain.

The key insight is that firms use layoffs not only to reduce labor input, but also to select which workers to keep. When layoffs disproportionately remove low-productivity matches, they improve the composition of the remaining workforce. This increases the value of retained workers and reduces firms’ incentives to cut wages. As a result, layoffs and wage rigidity are complements rather than substitutes: wage rigidity emerges precisely in states where firms engage in selection.

A central contribution of the paper is to show that this mechanism is both empirically identifiable and quantitatively important. It is, to my knowledge, the first endogenous mechanism that jointly accounts for wage rigidity, layoffs, and their heterogeneity across worker tenure. In particular, it explains why workers who are more exposed to layoffs—such as recent hires—also exhibit more rigid wages within the same firm, a pattern that is not accounted for by existing models.

To formalize this mechanism, I develop a model of a frictional labor market with directed search and dynamic contracts. Firms employ workers of heterogeneous match quality under asymmetric information: match productivity is privately observed by firms but not by workers. Firms post vacancies in submarkets offering promised lifetime utility and commit to long-term contracts that specify wages, layoff probabilities, and continuation values.

A key feature of the model is that contracts are cohort-specific. Workers hired under the same contract are treated identically and share the same wage path and layoff risk as long as they remain employed. Because match quality is unobserved and cannot be priced, firms cannot differentiate workers within a cohort through wages. Instead, layoffs are the only instrument through which firms can affect the composition of workers within a cohort.

This generates a tight link between selection and wage dynamics. Following negative shocks, firms lay off workers and disproportionately remove low-productivity matches. This improves the composition of the remaining workers within a cohort, increasing their expected productivity and the value of retaining them. As a result, cohorts that experience more intense selection exhibit stronger wage rigidity.

The model therefore delivers heterogeneity across cohorts in both layoffs and wage responses. Workers with lower continuation values—such as recent hires—are more exposed to layoffs and therefore experience stronger selection. In turn, their wages respond less to shocks. More senior workers, who face lower layoff risk, experience weaker selection and adjust primarily through wages. This within-firm, cross-cohort variation is central to the empirical analysis and provides a direct way to distinguish the selection mechanism from alternative explanations of wage rigidity. Importantly, because workers within a firm are exposed to the same productivity shocks, this cross-cohort variation isolates the role of selection rather than differences in firm-level conditions.

I test these predictions using French matched employer–employee data combined with firm-level production data. The empirical analysis proceeds in two steps. I first document a simple but important fact about the joint behavior of layoffs and wages across firms. I then use within-firm variation across worker tenure to identify the mechanism proposed in the model.

Across firms, I document a strong negative relationship between layoffs and wage pass-through. Firms that adjust employment more in response to shocks exhibit significantly more rigid wages. Following a negative productivity shock normalized to 100 percent, wages in low-layoff firms fall by 1.7 percent, whereas wages in high-layoff firms slightly increase. This pattern is robust and reflects a basic empirical regularity: firms that rely more on layoffs adjust wages less. This fact is difficult to reconcile with standard models of wage rigidity, in which wage dynamics are primarily driven by insurance motives and do not depend directly on firms’ layoff behavior. The evidence therefore suggests that mechanisms generating wage rigidity should account for how firms adjust employment.

I then turn to within-firm variation across worker tenure to isolate the selection mechanism. Workers within the same firm are exposed to identical productivity shocks, but differ in their continuation values and exposure to layoffs. I find that workers with less than two years of tenure face layoff rates between 5 and 11 percent, compared with around 2 percent for workers with four to five years of tenure. At the same time, wages of junior workers are substantially less responsive to firm-level shocks and even increase following negative shocks, while wages of more senior workers decline. These patterns mirror the model’s predictions: workers who are more exposed to layoffs experience stronger selection and therefore exhibit more rigid wages. I also show that wages respond positively to layoffs within a worker’s cohort and that layoffs increase wage dispersion across tenure groups, providing direct evidence for the selection channel.

To assess the quantitative importance of this mechanism, I calibrate the model to the French economy using simulated method of moments, targeting labor market flows, firm-level productivity dynamics, and wage moments. The model successfully reproduces the differential response of wages and layoffs across worker tenure, even though these patterns are not directly targeted. Following negative productivity shocks, the model generates large increases in layoff risk for junior workers but minimal wage declines, while more senior workers adjust primarily through wages.

France provides a particularly suitable setting for this analysis because of its strong labor market institutions, which allow me to compare the selection mechanism to alternative explanations for wage rigidity and layoffs. In particular, I incorporate two key policies: severance pay and the minimum wage.

I first consider severance pay, which represents a natural alternative explanation for the observed relationship between layoffs and tenure. In France, severance pay increases with tenure and is often more generous in practice than mandated by law. This could mechanically generate lower layoffs among senior workers. To address this, I allow firms to choose contract-specific severance payments in the model. I find that firms optimally offer substantial severance pay, increasing with tenure and exceeding legal requirements. Intuitively, firms use severance pay to smooth workers’ utility, making it less costly to separate from workers ex post rather than compensating them ex ante through higher wages. Importantly, while severance pay affects the level of layoffs, it does not generate the joint pattern of layoffs and wage rigidity across cohorts, which continues to be driven by selection.

I then turn to the role of the minimum wage as a source of exogenous wage rigidity. I first show that removing the minimum wage has only a limited effect on the heterogeneity of layoffs and wage pass-through across worker tenure. I then evaluate the impact of increasing the minimum wage on employment dynamics. In standard frameworks, restricting wage adjustments leads firms to substitute toward layoffs, amplifying employment fluctuations. In contrast, in this environment the effects are muted: a 20 percent increase in the minimum wage raises the annual layoff rate by only 0.2 percentage points and reduces the hiring rate by 1 percentage point.

The intuition is that layoffs already serve a valuable role as a selection device, even when wage adjustments are available. Moreover, firms optimally smooth wages through dynamic contracts, so the minimum wage constraint is often not strongly binding. When it does bind, firms adjust the timing of wages—frontloading compensation rather than raising it permanently—further limiting its impact on employment. Overall, these results suggest that when wage rigidity arises endogenously from selection, additional exogenous constraints have limited effects on equilibrium outcomes. This highlights that the effects of labor market policies depend critically on whether wage rigidity is imposed or arises endogenously from firms’ optimal decisions.

Taken together, these results suggest that selection is a key and previously missing component in understanding wage rigidity. More broadly, the paper highlights that wage dynamics cannot be understood without accounting for how firms choose which workers to retain, not just how they adjust pay.

\textbf{Related Literature}

This paper contributes to the literature on the origins of wage rigidity. Broadly, two approaches exist. One treats wage rigidity as an endogenous outcome of efficient contracting between firms and workers \parencite{barro1977,thomas1988,balke2022,elsby2024}. The other models rigid wages as the result of exogenously imposed frictions, such as nominal rigidities or institutional constraints \parencite{christiano2005,nekarda2020,blanco2025}. In the former, wage rigidity arises from insurance and incentive considerations; in the latter, it is imposed directly and often leads firms to adjust employment instead. This paper proposes a different mechanism: wage rigidity arises endogenously from firms’ selection of which workers to retain. By linking wage dynamics to layoffs, the paper provides a unified explanation for the joint behavior of wages and employment.

The paper is closest to the literature on dynamic wage contracts in frictional labor markets. Building on \textcite{holmstrom1982} and \textcite{thomas1988}, a number of papers study optimal contracts in environments with search frictions and on-the-job search, including \textcite{burdett2003}, \textcite{menzio2011}, \textcite{balke2022}, \textcite{souchier2022}, \textcite{cedomir2024}, and \textcite{elsby2024}. In these models, wage rigidity arises from the trade-off between insuring risk-averse workers and providing incentives to retain them. While they generate smooth wage dynamics, they do not account for the joint behavior of wages and layoffs across workers, nor for the systematic relationship between layoff risk and wage pass-through within firms. This paper introduces a distinct channel: layoffs alter workforce composition and therefore the value of retaining workers, which feeds back into wage-setting decisions. As a result, wage rigidity arises from selection rather than insurance.

The paper also relates to search-and-matching models with firm dynamics and heterogeneous workers, including \textcite{acemoglu2014}, \textcite{elsby2013}, \textcite{kaas2015}, \textcite{schaal2017}, \textcite{gulyas2020}, \textcite{bilal2022}, \textcite{elsby2022}, \textcite{mccrary2022}, and \textcite{rudanko2023}. These models study how firms adjust employment in response to shocks, often focusing on bargaining or wage posting. In contrast, this paper emphasizes the interaction between employment adjustment and wage dynamics under dynamic contracts. A key implication of the model is that heterogeneity in layoffs and wage responses arises across cohorts within the same firm, a feature that is not present in models where workers are treated symmetrically or wages are determined through static bargaining.

More broadly, the paper connects to work on match heterogeneity and endogenous separations, such as \textcite{berger2011}, \textcite{menzio2011}, and \textcite{gregory2021}. In these models, low-productivity matches are more likely to separate, generating cyclical patterns in employment and productivity. This paper builds on that insight but focuses on its implications for wage dynamics. By allowing firms to adjust both wages and employment, it shows that selection through layoffs affects not only which matches survive but also how surviving workers are paid.

On the empirical side, the paper relates to evidence on the wage–layoff trade-off and wage rigidity. \textcite{bewley1999} documents that firms are reluctant to cut wages and instead adjust employment. More recent work by \textcite{bertheau2025} and \textcite{davis2025} shows that layoffs are often used to target specific workers and that wage cuts are a poor substitute for separations. This paper provides a theoretical framework that rationalizes these findings: when firms value particular workers more than others, they optimally use layoffs for selection while smoothing wages for those they retain.

The paper also contributes to the empirical literature on wage pass-through and layoffs. Average wage pass-through estimates are consistent with \textcite{souchier2022} and cross-country evidence in \textcite{guvenen2017,guiso2020}. I document new evidence on heterogeneity in wage pass-through across worker tenure: junior workers experience lower or even negative pass-through, while more senior workers exhibit higher responsiveness. This pattern aligns with evidence on layoffs, such as \textcite{buhai2014}, who show that recently hired workers are more likely to be displaced. I also document a systematic relationship between layoffs and wage pass-through across firms, consistent with findings in \textcite{dias2013,enrlich2024}.

The remainder of the paper proceeds as follows. Section 2 presents the model. Section 3 provides empirical evidence. Section 4 presents the quantitative analysis.
